\documentclass[11pt]{article}
\usepackage{ideastyle}

\begin{document}

\ideatitle{StarFix}{Navigate without GPS: find your direction from a photo of the stars, like a spacecraft does.}

\section{The Idea}
Spacecraft use star trackers: cameras that recognize star patterns to work out
which way they are pointing. Sailors used the stars to find their position. StarFix
brings that to a phone. Take or upload a photo of the night sky; StarFix identifies
the stars, works out where the camera was pointing, and, combined with the time
and the phone's tilt, estimates where you are on Earth. Grok Voice walks you
through the result.

\section{Track Fit}
\begin{tabular}{@{}P{0.28\textwidth}P{0.66\textwidth}@{}}
\toprule
\req{Built with Cursor}{Image pipeline, star matching, and UI built in Cursor.}
\req{Grok Voice / Imagine}{\textbf{Voice} explains the fix; \textbf{Imagine} can produce a labeled sky chart.}
\req{Real space data}{Hipparcos or Yale Bright Star catalog; real star-tracker methods.}
\req{Navigation theme}{Navigation when GPS is unavailable: on Mars, in deep space, or when GPS is jammed.}
\req{Also eligible}{Big Red, Software; Hardware if using a Raspberry Pi camera.}
\bottomrule
\end{tabular}

\section{Features}
\subsection{MVP}
\begin{itemize}
  \item Upload a night-sky photo.
  \item Detect stars and match them to a catalog (``plate solving'').
  \item Show the photo with stars and constellations labeled, plus the pointing direction.
  \item Grok Voice explains: \emph{``You were facing northeast, about 40 degrees up.''}
\end{itemize}
\subsection{Stretch}
\begin{itemize}
  \item Estimate latitude and longitude from the star fix, time, and phone tilt.
  \item Live camera mode.
  \item Raspberry Pi camera as a DIY star tracker (Hardware track).
\end{itemize}

\section{Tech Stack and Data}
\begin{itemize}
  \item Star matching: astrometry.net (local install or its web API), or Python \texttt{tetra3}.
  \item Image processing: Python with OpenCV / Astropy.
  \item Frontend: React upload page with an annotated image.
  \item AI: Grok Voice API (and optionally Imagine).
\end{itemize}

\section{Limitations and Risks}
\begin{itemize}
\limitation{Hardest technically}{Reliable star matching is a serious algorithm. Writing it from scratch in 24 hours is very risky.}{Use an existing solver (astrometry.net, tetra3) and focus on the experience. Prototype this step first, Friday night.}
\limitation{Phone photos are hard}{Phones capture few stars, with noise, blur, and light pollution, especially in Ithaca in October (often cloudy).}{Use long-exposure or night-mode photos and keep a set of good sample images ready.}
\limitation{Daytime indoor judging}{You cannot take a sky photo during judging.}{Demo with pre-captured photos; show a short video of capturing one.}
\limitation{Location needs more than a photo}{A sky photo alone gives pointing direction, not position. Position also needs an accurate time and the camera's angle to the horizon.}{Keep location as a stretch goal; the MVP is ``which way am I pointing.''}
\limitation{Solver setup}{astrometry.net's local install needs large index files and is awkward on Windows; its web API can be slow (minutes per image).}{Test both Friday night; pre-solve demo images and cache the results.}
\limitation{Grok fit is thinner}{Voice explains results but isn't essential to the core function.}{Make the voice explanation educational: how sailors and spacecraft do this.}
\end{itemize}

\section{Scorecard}
\begin{tabular}{@{}P{0.25\textwidth}P{0.08\textwidth}P{0.6\textwidth}@{}}
\toprule
\rating{Demo-able}{3}{Pre-captured images only.}
\rating{Buildable in 24h}{2}{Depends heavily on the solver working.}
\rating{Navigation theme}{5}{Navigation without GPS.}
\rating{Wow factor}{5}{``Our phone found where it was from the stars.''}
\rating{Technical risk (5 = riskiest)}{5}{High.}
\bottomrule
\end{tabular}

\section{Verdict}
\textbf{High risk, high reward.} The most ``legendary'' concept, but only pick it
if a quick test Friday night shows the star solver working on a phone photo.

\end{document}
